%% kapa: exam=1 round=20 year=2009 period=2 form=A subject=appraisal_law subject_ko=부동산관계법규 seq=19 number=59
%% source: Q-Net 감정평가사 기출문제 (cst00309) artlSeq=5024178 file=extracted/1008/5024178/2009년 제20회 감정평가사 1차시험 기출문제/2교시.hwp
%% conversion: kordoc 4.24.51 -> tools/gen1.py (auto); review=proofread-sonnet; flags=none
\begin{question}{59}
부동산 가격공시 및 감정평가에 관한 법령상 주택가격의 공시에 관한 설명으로 옳지 \uline{않은} 것은?
\begin{choices}
  \item 국토해양부장관은 표준주택에 대하여 매년 공시기준일 현재의 적정가격을 조사․평가하고 중앙부동산평가위원회의 심의를 거쳐 이를 공시하여야 한다.
  \item 시장․군수 또는 구청장이 표준주택으로 선정된 단독주택에 대하여 개별주택가격을 결정․공시하지 아니할 경우에는 표준주택가격을 개별주택가격으로 본다.
  \item 표준주택가격은 국가ㆍ지방자치단체 등의 기관이 과세 등의 업무와 관련하여 주택의 가격을 산정하는 경우에 그 기준으로 활용된다.
  \item 국토해양부장관은 공동주택가격을 공시하기 위하여 공동주택의 가격을 산정한 때에는 공동주택소유자 그 밖의 이해관계인의 의견을 들어야 한다.
  \item 국토해양부장관은 공시기준일 이후에 토지의 분할․합병이나 건물의 신축 등이 발생한 경우에는 대통령령이 정하는 날을 기준으로 공동주택가격을 결정․공시하여야 한다.
\end{choices}
\end{question}
\answer{3}
